\chapter{Εισαγωγή}
\label{ch:introduction}

\section{Η Κβαντική Απειλή και η Post-Quantum Κρυπτογραφία}
\label{sec:quantum-threat}

% TODO: Αλγόριθμος Shor, harvest now decrypt later, CNSA 2.0

\section{Το Πλαίσιο Τυποποίησης του NIST και ο Αλγόριθμος \HQC}
\label{sec:nist-hqc}

% TODO: NIST rounds, τελικός υποψήφιος HQC

\section{Cryptographic Agility και Καθαρά Λογισμικές Επιταχύνσεις}
\label{sec:crypto-agility}

\subsection{Cryptographic Agility ως Απαίτηση}
\label{ssec:agility-requirement}

% TODO: SIKE example, ανωριμότητα PQC τοπίου

\subsection{Καθαρά Λογισμική Επιτάχυνση σε Ανοιχτή Αρχιτεκτονική}
\label{ssec:software-acceleration}

% TODO: Φιλοσοφία RISC-V, RVA23 profile

\section{Σκοπός και Συνεισφορά της Διατριβής}
\label{sec:contribution}

% TODO: Research gap, end-to-end βελτιστοποίηση στο RISC-V

\section{Δομή της Εργασίας}
\label{sec:structure}

Η εργασία οργανώνεται ως εξής:
\begin{itemize}
  \item Στο \autoref{ch:theoretical-background} παρουσιάζεται το θεωρητικό υπόβαθρο.
  \item Στο \autoref{ch:risc-v-architecture} αναλύεται η αρχιτεκτονική RISC-V.
  \item Στο \autoref{ch:related-work} ανασκοπούνται οι σχετικές εργασίες.
  \item Στο \autoref{ch:experimental-setup} περιγράφεται η πειραματική υποδομή.
  \item Στο \autoref{ch:polymul-optimization} αναπτύσσεται η βελτιστοποίηση του πολυωνυμικού πολλαπλασιασμού.
  \item Στο \autoref{ch:kernel-optimization} παρουσιάζονται οι βελτιστοποιήσεις Keccak και Decode.
  \item Στο \autoref{ch:results} παρατίθενται τα πειραματικά αποτελέσματα.
  \item Στο \autoref{ch:discussion} γίνεται η συζήτηση.
  \item Στο \autoref{ch:conclusions} εξάγονται τα συμπεράσματα.
\end{itemize}
